\section{Полёты БПЛА в России в 2026 году}

\begin{caution}[Важно помнить]
Правила меняются очень быстро — особенно региональные запреты. Ниже — состояние на сентябрь 2026 года.
Перед любым реальным полётом проверяйте актуальные документы: сайт Росавиации (\url{favt.gov.ru}),
сайт правительства своего региона и Госуслуги.
\end{caution}

\subsection{Основные термины (по Воздушному кодексу РФ)}
\begin{itemize}
\item \term{БВС} (беспилотное воздушное судно) — воздушное судно, управляемое внешним пилотом, который находится вне борта. В быту — «дрон», «БПЛА», «коптер».
\item \term{БАС} (беспилотная авиационная система) — комплекс: одно или несколько БВС + \term{станция внешнего пилота} (пульт, наземная станция) + \term{линии управления и контроля} (радиоканалы) + средства обеспечения (зарядка, транспорт и т.\,п.).
\item \term{Внешний пилот} — лицо, управляющее БВС с земли. Для БВС массой более 30~кг нужно \term{свидетельство внешнего пилота}.
\item \term{ИВП} — использование воздушного пространства. Полёт БВС на улице — это ИВП, для которого действуют \term{Федеральные правила ИВП} (Постановление Правительства РФ № 138 от 11.03.2010).
\item \term{VLOS / BVLOS} — полёт в пределах / за пределами прямой визуальной видимости пилота.
\end{itemize}

\subsection{Учёт и регистрация (федеральный уровень)}

\begin{table}[H]\centering\small
\renewcommand{\arraystretch}{1.3}
\begin{tabularx}{\linewidth}{@{}p{30mm} X X@{}}
\toprule
\textbf{Масса БВС} & \textbf{Что требуется} & \textbf{Документ} \\
\midrule
до \SI{150}{\gram} & Учёт не нужен. \emph{Но} региональные запреты и правила ИВП действуют! & ПП № 658 \\
\SIrange{0.15}{30}{\kilo\gram} & \term{Учёт} в Росавиации (через Госуслуги). Учётный номер наносится на корпус до первого полёта & ПП РФ № 658 от 25.05.2019 \\
более \SI{250}{\gram} & С \textbf{1 марта 2026} — обязательное оснащение аппаратурой удалённой идентификации / мониторинга и подключение к ГИС «ЭРА-ГЛОНАСС» (координаты, скорость, владелец — в реальном времени) & Постановление Правительства, Минтранс \\
более \SI{30}{\kilo\gram} & \term{Государственная регистрация}, сертификат лётной годности, свидетельство внешнего пилота & Воздушный кодекс РФ \\
\bottomrule
\end{tabularx}
\caption{Требования к БВС в зависимости от массы}
\end{table}

\subsection{Как законно выполнить полёт (общий порядок)}
\begin{enumerate}
\item Поставить БВС на учёт (если $m > \SI{150}{\gram}$), нанести учётный номер; с 2026 г. — подключить к «ЭРА-ГЛОНАСС» (если $m > \SI{250}{\gram}$).
\item Проверить, нет ли в месте полёта \term{запретных зон}, \term{зон ограничения полётов}, диспетчерских зон аэродромов (например, Пулково), зон вокруг режимных объектов. Проверить \textbf{региональный запрет}.
\item Над населённым пунктом — получить \term{разрешение органа местного самоуправления} (администрации).
\item Подать \term{план полёта} / заявку на установление местного режима в \term{зональный центр ЕС ОрВД} (через систему СППИ, обычно не позднее чем за сутки), в день полёта — связаться с диспетчером.
\item Выполнять полёт в светлое время суток, в пределах видимости, на безопасной высоте (типично до \SI{150}{\metre}), не над людьми и массовыми мероприятиями.
\end{enumerate}

\begin{keybox}[Полёты в помещении]
Полёт внутри здания (спортзал, аудитория, закрытая арена) \textbf{не является использованием воздушного пространства} РФ, поэтому план полёта и разрешения ЕС ОрВД не нужны. Именно поэтому учебные полёты и FPV-гонки в 2026 году проходят в залах. Нужно лишь согласие владельца помещения и соблюдение техники безопасности.
\end{keybox}

\subsection{Региональные запреты: откуда они берутся}
Указом Президента РФ № 757 от 19.10.2022 в субъектах РФ введены режимы готовности (базовый, средний, повышенный). На их основании губернаторы вправе вводить ограничения, в том числе \textbf{запрет на использование БПЛА} на территории региона. В 2025–2026 гг. такие запреты действуют в большинстве регионов европейской части России. Дополнительно при атаках БПЛА объявляется режим «\term{беспилотная опасность}»: могут ограничиваться мобильный интернет, закрываться аэропорты (режим «Ковёр»).

\subsection{Ленинградская область}
\begin{itemize}
\item С \textbf{мая 2023 года} действует \textbf{запрет на использование беспилотных воздушных судов} на всей территории области — постановление губернатора А.\,Ю.~Дрозденко.
\item Исключения — работы по специальному разрешению регионального \term{оперативного штаба} (обеспечение безопасности, инфраструктура, госзадачи). Для полётов тяжёлых БВС (свыше 30~кг) также требуется решение оперштаба.
\item Запрет действует до снятия режима базовой готовности, введённого Указом № 757.
\item В 2026 году неоднократно объявлялась «опасность БПЛА» (например, 28 июля 2026) — в это время возможны ограничения связи и транспорта.
\end{itemize}

\subsection{Санкт-Петербург}
\begin{itemize}
\item Запрет на запуск БПЛА продлён постановлением губернатора А.\,Д.~Беглова \textbf{до 31 декабря 2026 года}.
\item Тем же документом запрещена публикация фото/видео и сведений о последствиях атак БПЛА.
\item Рядом — аэропорт Пулково (диспетчерская зона), многочисленные запретные зоны над центром города.
\end{itemize}

\subsection{Ответственность}
\begin{table}[H]\centering\small
\renewcommand{\arraystretch}{1.25}
\begin{tabularx}{\linewidth}{@{}X r r r@{}}
\toprule
\textbf{Нарушение (КоАП РФ)} & \textbf{Граждане} & \textbf{Должностные} & \textbf{Юр. лица} \\
\midrule
ст. 11.4 ч.\,1 — нарушение Федеральных правил ИВП & \numrange{20}{50} тыс.\,руб. & \numrange{100}{150} тыс.\,руб. & \numrange{250}{300} тыс.\,руб. \\
ст. 11.4 — ИВП лицом, не имеющим на это права & \numrange{30}{50} тыс.\,руб. & — & — \\
Нарушение регионального запрета (режим готовности, ст. 20.6.1) & штраф & штраф & штраф \\
\bottomrule
\end{tabularx}
\caption{Штрафы за незаконные полёты (ориентировочно; суммы уточнять в действующей редакции КоАП)}
\end{table}
Кроме штрафа, дрон могут изъять, а силовые структуры вправе подавить или уничтожить БВС, нарушающий порядок ИВП. Полёт над охраняемыми объектами может повлечь уголовную ответственность.

\begin{exam}
\begin{enumerate}
\item Какие БВС подлежат учёту, а какие — государственной регистрации?
\item Что изменилось с 1 марта 2026 года?
\item Почему в Ленинградской области нельзя просто так запустить дрон? Какой документ это запрещает?
\item Нужен ли план полёта для полёта в спортзале? Почему?
\end{enumerate}
\end{exam}
